% GENERATED FILE. Edit canonical lessons, then run npm run generate:books.
% canonical-source-hash: fnv1a64:c45da731f18e679d
% canonical-lessons: BN-C178-dater-daktar, BN-C178-domkol, BN-C178-dhatri, BN-C178-bikreta, BN-C178-kreta, BN-R178-first-pass-words-from-chapters-175-176, BN-R178-second-pass-words-from-chapters-177-178

\chapter{Dentist, Fire brigade, Midwife, Seller, Buyer}
\label{ch:bn-dentist-fire-brigade-midwife-seller-buyer}
\hlchaptermodality{178}

\begin{chapteropening}
\textbf{By the end of this chapter:} \emph{I can say a dentist, the fire brigade, a midwife, a seller and a buyer.}

Complete the last lesson of chapter 178: I can say a dentist, the fire brigade, a midwife, a seller and a buyer.
\end{chapteropening}

\section[\texorpdfstring{\bn{দাঁতের} \bn{ডাক্তার} (dā̃ter ḍāktār)}{দাঁতের ডাক্তার (dā̃ter ḍāktār)}]{\bn{দাঁতের} \bn{ডাক্তার} (dā̃ter ḍāktār) — a dentist}

\label{lesson:BN-C178-dater-daktar}



\begin{quote}
Before the new one: say the Bengali for a scientist, then the Bengali for a singer.
\end{quote}

\subsection*{You'll want to know: \bn{দাঁতের} \bn{ডাক্তার}}
\textbf{\bn{দাঁতের} \bn{ডাক্তার}} — \emph{dā̃ter ḍāktār} — \textquotedblleft{}a dentist\textquotedblright{}.

\begin{grammarlens}[title={What you've built}]
One more word for this chapter.
\end{grammarlens}

\subsection*{Your turn}
\begin{itemize}
\raggedright
  \item \emph{Say it:} \emph{dā̃ter ḍāktār}
  \item \emph{Say it:} \emph{dā̃ter ḍāktār}, once more
  \item \emph{Say it:} say \emph{dā̃ter ḍāktār}
  \item \emph{Recall:} say the Bengali for a scientist, then the Bengali for a singer, then say \emph{dā̃ter ḍāktār} again
\end{itemize}

\begin{tcolorbox}[breakable,colback=teal!4,colframe=teal!35!black,title={Before you move on}]
What does \bn{দাঁতের} \bn{ডাক্তার} mean? (\textquotedblleft{}A dentist\textquotedblright{}.) Say it once more.
\end{tcolorbox}
\section[\texorpdfstring{\bn{দমকল} (dômkôl)}{দমকল (dômkôl)}]{\bn{দমকল} (dômkôl) — the fire brigade}

\label{lesson:BN-C178-domkol}



\begin{quote}
Before the new one: say the Bengali for a singer, then the Bengali for a dentist.
\end{quote}

\subsection*{You'll want to know: \bn{দমকল}}
\textbf{\bn{দমকল}} — \emph{dômkôl} — \textquotedblleft{}the fire brigade\textquotedblright{}.

\begin{grammarlens}[title={What you've built}]
One more word for this chapter.
\end{grammarlens}

\subsection*{Your turn}
\begin{itemize}
\raggedright
  \item \emph{Say it:} \emph{dômkôl}
  \item \emph{Say it:} \emph{dômkôl}, once more
  \item \emph{Say it:} say \emph{dômkôl}
  \item \emph{Recall:} say the Bengali for a singer, then the Bengali for a dentist, then say \emph{dômkôl} again
\end{itemize}

\begin{tcolorbox}[breakable,colback=teal!4,colframe=teal!35!black,title={Before you move on}]
What does \bn{দমকল} mean? (\textquotedblleft{}The fire brigade\textquotedblright{}.) Say it once more.
\end{tcolorbox}
\section[\texorpdfstring{\bn{ধাত্রী} (dhātrī)}{ধাত্রী (dhātrī)}]{\bn{ধাত্রী} (dhātrī) — a midwife}

\label{lesson:BN-C178-dhatri}



\begin{quote}
Before the new one: say the Bengali for a dentist, then the Bengali for the fire brigade.
\end{quote}

\subsection*{You'll want to know: \bn{ধাত্রী}}
\textbf{\bn{ধাত্রী}} — \emph{dhātrī} — \textquotedblleft{}a midwife\textquotedblright{}.

\begin{grammarlens}[title={What you've built}]
One more word for this chapter.
\end{grammarlens}

\subsection*{Your turn}
\begin{itemize}
\raggedright
  \item \emph{Say it:} \emph{dhātrī}
  \item \emph{Say it:} \emph{dhātrī}, once more
  \item \emph{Say it:} say \emph{dhātrī}
  \item \emph{Recall:} say the Bengali for a dentist, then the Bengali for the fire brigade, then say \emph{dhātrī} again
\end{itemize}

\begin{tcolorbox}[breakable,colback=teal!4,colframe=teal!35!black,title={Before you move on}]
What does \bn{ধাত্রী} mean? (\textquotedblleft{}A midwife\textquotedblright{}.) Say it once more.
\end{tcolorbox}
\section[\texorpdfstring{\bn{বিক্রেতা} (bikretā)}{বিক্রেতা (bikretā)}]{\bn{বিক্রেতা} (bikretā) — a seller}

\label{lesson:BN-C178-bikreta}



\begin{quote}
Before the new one: say the Bengali for the fire brigade, then the Bengali for a midwife.
\end{quote}

\subsection*{You'll want to know: \bn{বিক্রেতা}}
\textbf{\bn{বিক্রেতা}} — \emph{bikretā} — \textquotedblleft{}a seller\textquotedblright{}.

\begin{grammarlens}[title={What you've built}]
One more word for this chapter.
\end{grammarlens}

\subsection*{Your turn}
\begin{itemize}
\raggedright
  \item \emph{Say it:} \emph{bikretā}
  \item \emph{Say it:} \emph{bikretā}, once more
  \item \emph{Say it:} say \emph{bikretā}
  \item \emph{Recall:} say the Bengali for the fire brigade, then the Bengali for a midwife, then say \emph{bikretā} again
\end{itemize}

\begin{tcolorbox}[breakable,colback=teal!4,colframe=teal!35!black,title={Before you move on}]
What does \bn{বিক্রেতা} mean? (\textquotedblleft{}A seller\textquotedblright{}.) Say it once more.
\end{tcolorbox}
\section[\texorpdfstring{\bn{ক্রেতা} (kretā)}{ক্রেতা (kretā)}]{\bn{ক্রেতা} (kretā) — a buyer}

\label{lesson:BN-C178-kreta}



\begin{quote}
Before the new one: say the Bengali for a midwife, then the Bengali for a seller.
\end{quote}

\subsection*{You'll want to know: \bn{ক্রেতা}}
\textbf{\bn{ক্রেতা}} — \emph{kretā} — \textquotedblleft{}a buyer\textquotedblright{}.

\begin{grammarlens}[title={What you've built}]
One more word for this chapter.
\end{grammarlens}

\subsection*{Your turn}
\begin{itemize}
\raggedright
  \item \emph{Say it:} \emph{kretā}
  \item \emph{Say it:} \emph{kretā}, once more
  \item \emph{Say it:} say \emph{kretā}
  \item \emph{Recall:} say the Bengali for a midwife, then the Bengali for a seller, then say \emph{kretā} again
\end{itemize}

\begin{tcolorbox}[breakable,colback=teal!4,colframe=teal!35!black,title={Before you move on}]
What does \bn{ক্রেতা} mean? (\textquotedblleft{}A buyer\textquotedblright{}.) Say it once more.
\end{tcolorbox}
\section[(dialogue)]{Words from chapters 175-176 — a first pass}

\label{lesson:BN-R178-first-pass-words-from-chapters-175-176}



\begin{quote}
Say the words for a seller and a buyer.
\end{quote}

\subsection*{Your turn}
Say these aloud:

\begin{itemize}
\raggedright
  \item a stench, a duster, a bonti, a grinding stone, gas
  \item a profession, a debt, a bill, an expense, savings
\end{itemize}

\begin{tcolorbox}[breakable,colback=teal!4,colframe=teal!35!black,title={Before you move on}]
A stench? (\textbf{\bn{দুর্গন্ধ}}, \emph{durgôndhô}.) A duster? (\textbf{\bn{ঝাড়ন}}, \emph{jhāṛôn}.) A bonti? (\textbf{\bn{বঁটি}}, \emph{bõṭi}.) A grinding stone? (\textbf{\bn{শিলনোড়া}}, \emph{shilnoṛā}.) Gas? (\textbf{\bn{গ্যাস}}, \emph{gæs}.) A profession? (\textbf{\bn{পেশা}}, \emph{peshā}.) A debt? (\textbf{\bn{ধার}}, \emph{dhār}.) A bill? (\textbf{\bn{বিল}}, \emph{bil}.) An expense? (\textbf{\bn{খরচ}}, \emph{khôrôch}.) Savings? (\textbf{\bn{সঞ্চয়}}, \emph{sôñchôy}.)
\end{tcolorbox}
\section[(dialogue)]{Words from chapters 177-178 — a second pass}

\label{lesson:BN-R178-second-pass-words-from-chapters-177-178}



\begin{quote}
Say the words for a judge and a student.
\end{quote}

\subsection*{Your turn}
Say these aloud:

\begin{itemize}
\raggedright
  \item a judge, a student, an engineer, a scientist, a singer
  \item a dentist, the fire brigade, a midwife, a seller, a buyer
\end{itemize}

\begin{tcolorbox}[breakable,colback=teal!4,colframe=teal!35!black,title={Before you move on}]
A judge? (\textbf{\bn{বিচারক}}, \emph{bichārôk}.) A student? (\textbf{\bn{ছাত্র}}, \emph{chhātrô}.) An engineer? (\textbf{\bn{ইঞ্জিনিয়ার}}, \emph{injiniyār}.) A scientist? (\textbf{\bn{বিজ্ঞানী}}, \emph{bigyānī}.) A singer? (\textbf{\bn{গায়ক}}, \emph{gāyôk}.) A dentist? (\textbf{\bn{দাঁতের} \bn{ডাক্তার}}, \emph{dā̃ter ḍāktār}.) The fire brigade? (\textbf{\bn{দমকল}}, \emph{dômkôl}.) A midwife? (\textbf{\bn{ধাত্রী}}, \emph{dhātrī}.) A seller? (\textbf{\bn{বিক্রেতা}}, \emph{bikretā}.) A buyer? (\textbf{\bn{ক্রেতা}}, \emph{kretā}.)
\end{tcolorbox}
